% HD-G13. GPT-6 Astra Pro, /root/g34. 21 November 2026.
% Reasoning effort inherited; numerical setting not separately exposed.
\section{The exact trace obstruction for a genus-four Abel ideal}
\label{sec:genus4-trace}

\paragraph{HD-G13.}
Let $C$ be a smooth connected complex curve of genus four, let
$i:C\hookrightarrow A=J(C)$ be an Abel embedding, and put
$I=I_{i(C)}$. Set
\[
 W=H^0(C,K_C),\qquad
 \mu:\operatorname{Sym}^2W\longrightarrow H^0(C,K_C^2),
 \qquad I_2(C)=\ker\mu.
\]
For the full degree-two semiregularity trace
\[
 \sigma_I^2:\operatorname{Ext}^2_A(I,I)
   \longrightarrow
   H^2(A,\mathcal O_A)\oplus H^3(A,\Omega_A^1)
                         \oplus H^4(A,\Omega_A^2),
\]
there is an explicit isomorphism
\[
 \boxed{I_2(C)^*\ \xrightarrow{\ \sim\ }\ker\sigma_I^2.}
                                                               \tag{G13.1}
\]
It is induced by the Hochschild evaluation of polarized
infinitesimal deformations, as specified in (G13.9) below.
The kernel has dimension one for a nonhyperelliptic genus-four
curve and dimension three for a hyperelliptic genus-four curve.
In particular this particular ideal is not semiregular.

Moreover the stabilizer of $I$ inside the actual
translation/tensor group $A\times\widehat A$ is trivial.
Thus no nontrivial finite subgroup of that group can give this
ideal an equivariant structure, and the weak equivariant
semiregularity allowed in Perry's Theorem~1.2 does not remove
its trace kernel. This is a statement about the specified ideal
and the explicitly described constructions below.

\paragraph{The normal sequence and its exact cohomology maps.}
Write $V=H^0(A,T_A)=W^*$. Restriction along the Abel embedding
identifies $H^1(A,\mathcal O_A)$ with $H^1(C,\mathcal O_C)=W^*$.
The differential of $i$ gives
\[
 0\longrightarrow T_C\longrightarrow V\otimes\mathcal O_C
          \longrightarrow N\longrightarrow0,
 \qquad N=N_{C/A}.                                      \tag{G13.2}
\]
The dual of $H^1(T_C)\to V\otimes H^1(\mathcal O_C)$ is
the complete multiplication map
\[
 W\otimes W\longrightarrow H^0(K_C^2).
\]
Indeed the dual of the differential in (G13.2) is the evaluation
map $W\otimes\mathcal O_C\to K_C$, and Serre duality produces
precisely that multiplication map. Put $h=\operatorname{rank}\mu$.
Since $h^1(T_C)=9$, the long exact sequence gives
\[
 h^0(N)=13-h,\qquad h^1(N)=16-h,
\]
and the actual quotient map is
\[
 V\otimes H^1(\mathcal O_C)=W^*\otimes W^*
 \longrightarrow H^1(N)
  =(W^*\otimes W^*)/\operatorname{im}(\mu^*).             \tag{G13.3}
\]
Here $\operatorname{im}(\mu^*)$ lies in
$\operatorname{Sym}^2W^*$ and equals the annihilator of $I_2(C)$.

Choose a nonzero invariant volume form on $A$. The determinant
of (G13.2) then identifies $\det N$ with $K_C$, hence
$\bigwedge^2N$ with $N^*\otimes K_C$ by
$(n_1\wedge n_2)(n_3)=n_1\wedge n_2\wedge n_3$.
Dualizing (G13.2) and tensoring by $K_C$ gives
\[
 H^0(\textstyle\bigwedge^2N)
   =\ker\{W\otimes W\longrightarrow H^0(K_C^2)\}
   =\bigwedge^2W\oplus I_2(C).                           \tag{G13.4}
\]
Its dimension is $16-h$. The ambient map
$\bigwedge^2V\to H^0(\bigwedge^2N)$ is injective and its image
is the first summand in (G13.4). To see the map itself, take
dual bases $e_1,\ldots,e_4$ of $V$ and $x_1,\ldots,x_4$ of $W$,
and volume $x_1\wedge\cdots\wedge x_4$. Then $e_1\wedge e_2$
maps, through the preceding determinant identification, to
$x_4\otimes x_3-x_3\otimes x_4$; the other five vectors map to
the corresponding signed alternating complementary pairs.
These six images form a basis of $\bigwedge^2W$.

\paragraph{All self-Ext dimensions, with no smoothness assumption
on the Hilbert scheme.}
The smooth curve is a regular immersion of codimension three.
Its local Koszul resolution gives
\[
 \mathcal H om(I,I)=\mathcal O_A,\qquad
 \mathcal E xt^1(I,I)=i_*N,\qquad
 \mathcal E xt^2(I,I)=i_*\bigwedge^2N,
 \qquad\mathcal E xt^{>2}(I,I)=0.                         \tag{G13.5}
\]
For clarity, the two higher identifications follow by applying
$R\mathcal H om(\mathcal O_C,-)$ to
$0\to I\to\mathcal O_A\to\mathcal O_C\to0$.
The only nonzero $\mathcal E xt^j(\mathcal O_C,\mathcal O_A)$
is $\det N$ in degree three; its map to
$\mathcal E xt^3(\mathcal O_C,\mathcal O_C)=\det N$ is the
identity in Koszul coordinates. This gives exactly (G13.5).

Applying $R\operatorname{Hom}(I,-)$ to the same sequence yields
\[
 0\longrightarrow H^0(N)\longrightarrow
  \operatorname{Ext}^1(I,I)\longrightarrow H^1(\mathcal O_A)
  \longrightarrow0.
\]
The last arrow splits by the scalar action on $I$: its trace
is multiplication by the rank, which is one. Equivalently,
$\operatorname{Ext}^1(I,\mathcal O_A)=H^1(\mathcal O_A)$
and the composition of the scalar action with $I\to\mathcal O_A$
is the identity on this group. Therefore
\[
 \operatorname{ext}^1(I,I)=17-h.
\]

Let $\theta$ be the principal polarization. The Abel embedding
has $[C]=\theta^3/6$: in an integral symplectic basis of
$H^1(A)$, restriction to $C$ pairs the $a_j,b_j$ with
$\int_Ca_jb_l=\delta_{jl}$, and its Poincare-dual class is the
sum of the four products omitting one pair, namely $\theta^3/6$.
Grothendieck--Riemann--Roch for $i$ and $\chi(\mathcal O_C)=-3$
give
\[
 \operatorname{ch}(I)=1-\theta^3/6+3[\mathrm{pt}],\qquad
 \operatorname{ch}(I^\vee)=1+\theta^3/6+3[\mathrm{pt}].
\]
The tangent bundle of $A$ is trivial, so Riemann--Roch gives
$\chi(I,I)=6$. Serre duality on the fourfold and
$\operatorname{hom}(I,I)=1$ now imply
\[
 \operatorname{ext}^2(I,I)=38-2h.                         \tag{G13.6}
\]
The three entries of total degree two in the local-to-global
Ext spectral sequence have dimensions
\[
 h^2(\mathcal O_A)=6,\qquad h^1(N)=16-h,
 \qquad h^0(\textstyle\bigwedge^2N)=16-h.
\]
Their sum is exactly (G13.6). Thus every one survives, giving
a filtration of $\operatorname{Ext}^2(I,I)$ with these three
associated graded spaces. No assertion of unobstructed embedded
curve deformations is involved.

\paragraph{The evaluation map has exact rank $28-h$.}
Let $\operatorname{ob}_I$ be Hochschild evaluation on $I$.
In the Atiyah--HKR description its domain is
\[
 HH^2(A)=H^2(\mathcal O_A)\oplus H^1(T_A)
                         \oplus H^0(\textstyle\bigwedge^2T_A).
\]
It respects the preceding filtration. Its three associated
graded maps are respectively
\[
 \begin{array}{ccl}
 H^2(\mathcal O_A)&\longrightarrow&H^2(\mathcal O_A),
       \quad\hbox{the identity},\\
 H^1(T_A)&\longrightarrow&H^1(N),
       \quad\hbox{the quotient in (G13.3)},\\
 H^0(\bigwedge^2T_A)&\longrightarrow&H^0(\bigwedge^2N),
       \quad\hbox{the ambient wedge map in (G13.4)}.
 \end{array}                                             \tag{G13.7}
\]
These are the actual leading maps of the Atiyah action. They
can be checked on the truncated Koszul resolution of the ideal
$(z_1,z_2,z_3)$: differentiating the differential in the direction
$\partial_{z_j}$ gives contraction by $e_j$ on the free exterior
resolution. Composing two such contractions sends
$e_1\wedge e_2\wedge e_3$ to the signed remaining generator.
After augmentation this is the corresponding class of $z_j$
in $(I/I^2)\otimes\det N=\bigwedge^2N$. The divided Atiyah
square supplies precisely the alternating wedge map. Changes
of local regular parameters respect these normal-bundle maps.

The first two components fill the first two filtration steps,
of total dimension $6+(16-h)$. The six bivectors have independent
images in the final quotient. Consequently
\[
 \operatorname{rank}\operatorname{ob}_I=28-h.              \tag{G13.8}
\]

\paragraph{Adjointness computes the full trace kernel.}
The same invariant volume form gives a nondegenerate Serre
pairing on $\operatorname{Ext}^2(I,I)$ and a nondegenerate
pairing between $HH^2(A)$ and $HH_{-2}(A)$. On the summand
$H^{2-p}(\bigwedge^pT_A)$, the latter is the Serre pairing
with $H^{p+2}(\Omega_A^p)$ after contraction with the volume.
The trace formula
$\operatorname{tr}(\xi\operatorname{At}_I^p/p!)$ shows directly
that $\sigma_I^2$ is adjoint to $\operatorname{ob}_I$, with the
signed Atiyah convention used in both maps. Indeed contraction
of the $p$th trace component against the polyvector is exactly
the Serre trace of $\xi$ composed with its Atiyah evaluation.
All the pairings are perfect; hence (G13.8) gives
\[
 \operatorname{rank}\sigma_I^2=28-h,
 \qquad\dim\ker\sigma_I^2=10-h=\dim I_2(C).
\]

Here is the promised identification of the kernel, rather than
only its dimension. The polarized tangent space inside
$H^1(T_A)=W^*\otimes W^*$ is $\operatorname{Sym}^2W^*$.
In coordinates $\theta=\sum x_jy_j$, contraction of
$\sum\kappa_{jl}\partial_{x_j}\otimes y_l$ with $\theta$
is zero exactly when the matrix $\kappa$ is symmetric.
The trace--evaluation identity gives
\[
 \sigma_I^2(\operatorname{ob}_I(\kappa))
       =\kappa\mathbin{\lrcorner}\operatorname{ch}(I)=0
       \qquad(\kappa\in\operatorname{Sym}^2W^*).
\]
The cubic polarization term contracts to
$\theta^2(\kappa\lrcorner\theta)/2=0$, while the top-degree
term would land in antiholomorphic degree five and vanishes.
By (G13.3), a symmetric $\kappa$ has nonzero image in $H^1(N)$
exactly when its functional on $I_2(C)$ is nonzero.
For $\kappa\in\operatorname{Ann}(I_2(C))$, its evaluation lies
in $H^2(\mathcal O_A)$; the trace on that filtration step is
injective, so its evaluation is zero. Thus evaluation descends
to the actual isomorphism
\[
 \operatorname{Sym}^2W^*/\operatorname{Ann}(I_2(C))
   =I_2(C)^*
 \ \xrightarrow{\ [\kappa]\mapsto\operatorname{ob}_I(\kappa)\ }
   \ker\sigma_I^2,                                      \tag{G13.9}
\]
whose injectivity follows already from its normal-bundle image.
The equal dimensions prove surjectivity.

For nonhyperelliptic $C$, the canonical embedding lies in
$\mathbb P^3$ and has degree six. There is at least one quadratic
equation since $\dim\operatorname{Sym}^2W=10>9=h^0(K_C^2)$.
There cannot be two independent quadrics: two without a common
component have intersection of degree four; a common plane
component would force the irreducible nondegenerate canonical
curve into that plane or into a line. Both alternatives contradict
the canonical embedding. Hence $h=9$ and $\dim I_2(C)=1$.
For hyperelliptic $C$, in a model $y^2=f(x)$ the canonical basis
is $x^jdx/y$ for $0\le j\le3$. Its products span precisely the
seven monomials $x^j(dx/y)^2$ for $0\le j\le6$.
Thus $h=7$ and $\dim I_2(C)=3$. The resulting exact dimensions are
\[
\begin{array}{c|rrrr}
 &\operatorname{ext}^1&\operatorname{ext}^2&
       \operatorname{rank}\sigma^2&\dim\ker\sigma^2\\\hline
 \text{nonhyperelliptic}&8&20&19&1\\
 \text{hyperelliptic}&10&24&21&3.
\end{array}
\]

\paragraph{Why an allowed finite action does not repair this ideal.}
Suppose a point $(a,P)\in A\times\widehat A$ gives an
isomorphism $t_a^*I\otimes P\simeq I$. Since $\det I=\mathcal O_A$
and $\operatorname{rank}I=1$, determinants give $P=\mathcal O_A$.
The isomorphism of ideals extends to their double duals, which
are $\mathcal O_A$; it is multiplication by a nonzero constant.
It therefore identifies the actual closed subschemes and gives
$t_a(C)=C$. The induced automorphism $\gamma$ of $C$ acts
identically on $H^0(K_C)$: the pullback
$i^*:H^0(\Omega_A^1)\to H^0(K_C)$ is an isomorphism and all
invariant forms on $A$ are translation invariant.

For a nonhyperelliptic curve the canonical embedding implies
$\gamma=1$. For a hyperelliptic curve the canonical morphism
shows that $\gamma$ is either the identity or the hyperelliptic
involution; the latter acts as $-1$ on the basis $x^jdx/y$ and
so is excluded. Thus $\gamma=1$ and then $a=0$.
The actual stabilizer is trivial. In particular there is no
nontrivial coherent action of a finite subgroup of
$A\times\widehat A$ on this ideal; taking that subgroup to be
trivial leaves exactly the trace kernel (G13.1).

\paragraph{The obstruction through the natural descent and square.}
Let $p:A'\to A$ be an actual isogeny and equip $p^*I$ with its
canonical coherent action by $H=\ker p$. Finite etale descent
identifies the invariant category with $D^{\mathrm{perf}}(A)$
and the descended object with $I$. Hence weak $H$-semiregularity
fails by (G13.1), not merely its stronger ambient-trace version.
This also exhausts the allowable point subgroups for this
pulled-back ideal. Its determinant is trivial, so a stabilizing
translation/tensor point again has trivial tensor coordinate.
A translation preserving $p^{-1}C$ gives $p(a)+C=C$ after
applying $p$, and the stabilizer just proved implies $p(a)=0$.
Conversely every element of $H$ preserves that inverse image.
Thus the full stabilizer of $p^*I$ is precisely $H\times\{0\}$.
For a subgroup of $H$, the descended object is a further finite
etale pullback of $I$. Pullback on self-Ext is injective by the
degree-normalized trace splitting, and trace naturality carries
the nonzero kernel (G13.1) into its trace kernel. Tensoring a
descent by a character line bundle does not change that kernel.
There is no additional coherence exception here. The ideal
$p^*I$ has scalar automorphisms, so two linearizations for the
canonical translation action differ by a character. More
generally, a coherence differing from the canonical translation
coherence by a scalar two-cocycle admits a linearization on this
ideal only when that cocycle is the coboundary of the scalars
comparing it with the canonical linearization. Changing those
comparisons identifies the action with the canonical one.
The scalar assertion for natural automorphisms of a translation
functor is $HH^0(A')=H^0(\mathcal O_{A'})=\mathbb C$; its negative
Hochschild groups vanish, so there are no further higher
automorphisms to alter this argument.

On a product of two such fourfolds the class
$\xi\boxtimes\operatorname{id}$, with $0\ne\xi\in\ker\sigma_I^2$,
is a nonzero degree-two self-Ext class. The external-product
trace is
$\sigma_I^2(\xi)\boxtimes\operatorname{ch}(I)=0$.
The same argument works for two different Abel ideals.
The explicit Fourier equivalences, line-bundle twists and
coordinate isomorphisms used in the secant-square construction
transport equivariant categories and their traces; duality and
shifts identify the same kernels by signed transpose.
Consequently this direct genus-four replacement of the natural
ideal input does not meet the permitted deformation theorem's
semiregularity hypothesis. This calculation makes no assertion
about ideal sheaves of surfaces, other complexes, or other CM
source constructions.
\hfill$\square$

\paragraph{Primary-source and verification scope.}
The characteristic Hochschild action through the divided Atiyah
powers is Buchweitz--Flenner, \emph{The global decomposition
theorem for Hochschild (co-)homology of singular spaces via the
Atiyah--Chern character}, Adv. Math. 217 (2008), 243--281,
Corollary~5.2.3. The trace definition is Section~6.1.2, and its
composition with evaluation is Theorem~6.4.2. These were reopened
in the original author version \url{https://arxiv.org/pdf/math/0606730},
printed pp.27, 29 and 31, respectively. The adjointness used
above is also proved directly from that trace formula and Serre
duality, and the local associated graded maps are computed from
the actual ideal Koszul resolution.
The equivariant condition is Perry, Definition~2.6 and
Lemma~2.7, printed pp.5--6, with trace functoriality in
Remark~2.3, p.5, and Theorem~1.2, p.3, in
\url{https://arxiv.org/pdf/2604.00511}. No deformation theorem
of Markman's manuscript is used in this obstruction calculation.
